\documentclass[11pt,a4paper]{article}
\usepackage[utf8]{inputenc}
\usepackage[T1]{fontenc}
\usepackage{amsmath,amssymb,amsfonts}
\usepackage{booktabs}
\usepackage{graphicx}
\usepackage{multirow}
\usepackage{hyperref}
\usepackage{xcolor}
\usepackage[margin=1in]{geometry}
\usepackage{natbib}
\usepackage{array}
\usepackage{tabularx}

\newcommand{\lzstar}{$\ell z^*$}
\newcommand{\lzm}{$\ell z_m$}

\title{IRT-FIT: Contamination Detection at the Minimal-Information Operating Point\\via Item Response Theory Person-Fit Statistics}

\author{Anonymous}
\date{}

\begin{document}
\maketitle

% ============================================================
% ABSTRACT
% ============================================================
\begin{abstract}
Benchmark contamination undermines the reliability of LLM evaluation, yet every existing detection method requires privileged access---training data, token probabilities, log-probs, or model weights.
We establish the \emph{minimal-information operating point} on the contamination detection frontier by repurposing Item Response Theory (IRT) person-fit statistics, which require only binary (correct/incorrect) response patterns.
Our framework, IRT-FIT, computes the \lzstar{} person-fit statistic under a 4-parameter logistic (4PL) IRT model and detects contamination as a shift in \lzstar{} relative to a clean baseline.
In controlled contamination experiments across four model families (Qwen2.5-7B, Llama-3.1-8B, Gemma-4-12B, Mistral-7B-v0.3) and three benchmarks, IRT-FIT achieves $|d| > 2.5$ (Cohen's $d$) at 15\% contamination on MMLU for all families, with monotonic dose--response under $\theta$-fixed analysis.
Critically, after GRPO reinforcement learning, IRT-FIT retains an AUROC of 0.890 ($\Delta = -0.035$), while Min-K\%++ and PPL-based methods collapse to near-random ($\Delta > -0.39$)---binary response patterns are structurally immune to RL-induced distributional shifts.
Against five baselines, IRT-FIT is competitive with token-probability methods (AUROC gap $= -0.035$ at 15\%, 95\% CIs overlapping) while requiring strictly less information.
Detection power follows a predictable function of benchmark size ($R^2 = 0.965$ across five benchmarks from 1.3K to 14K items), providing practitioners with actionable deployment guidelines.
\end{abstract}

% ============================================================
% 1. INTRODUCTION
% ============================================================
\section{Introduction}

Benchmark contamination---the inclusion of evaluation items in training data---is a growing threat to LLM evaluation integrity.
As training corpora scale to trillions of tokens and benchmarks proliferate, the probability of overlap increases, potentially inflating reported capabilities and misleading downstream decisions.

Existing detection methods span a hierarchy of information requirements.
N-gram overlap methods require access to training data.
Membership Inference Attacks (MIA) such as Min-K\%++ \citep{shi2024detecting} require per-token probability access.
PPL-based methods require sequence-level log-probabilities.
ConStat \citep{dekoninck2024constat} requires a second inference pass with rephrased items.
Time Travel probing \citep{oren2024proving} requires specially crafted temporal prompts.
Critically, Wang et al.\ (2025) demonstrate that \emph{all} ten log-probability-based methods they evaluate degrade to near-random performance on RL-trained reasoning models, because RL fundamentally alters token-level probability distributions while preserving which items are answered correctly.

A contaminated LLM is functionally analogous to an examinee with item preknowledge---a problem studied for decades in psychometrics under the framework of person-fit analysis.
Person-fit statistics such as $\ell z^*$ \citep{snijders2001} detect aberrant response patterns: examinees who answer hard items correctly but easy items incorrectly, inconsistent with their estimated ability level.
Crucially, these statistics operate on \emph{binary response patterns alone}---no token probabilities, no log-probs, no model weights.

We introduce IRT-FIT, a contamination detection framework that adapts IRT person-fit statistics to the LLM evaluation setting.
Our contributions are threefold:

\begin{enumerate}
\item \textbf{Minimal-information detection framework.}
We establish the minimal-information operating point on the contamination detection frontier.
IRT-FIT requires only binary response patterns (correct/incorrect per item) and achieves competitive detection performance against methods with strictly more information access (AUROC $= 0.925$ at 15\% contamination on MMLU, within 0.035 of Min-K\%++, with overlapping 95\% CIs).
After RL training, IRT-FIT retains $\text{AUROC} = 0.890$ while log-prob methods collapse.

\item \textbf{Detection characterization.}
Detection power is a predictable function of benchmark size and contamination level ($R^2 = 0.965$ across five benchmarks spanning 1.3K--14K items).
We identify and characterize the $\theta$-shift boundary condition: when ability is re-estimated on contaminated data, detection sensitivity can degrade, but fixing $\theta$ to the pre-contamination baseline restores monotonic dose--response across all model families.
The memorization fraction (0.33--0.38 across families at 15\% contamination) quantifies the gap between theoretical upper-bound power and observed detection.

\item \textbf{Empirical findings.}
LLM response patterns exhibit universal 4PL model misfit ($|lz^*| \gg 5$ in 57/60 model--benchmark cells), revealing that LLMs fundamentally differ from classical test-takers.
Domain-level analysis uncovers distinct contamination typologies (uniform, selective, domain-uneven) across models.
Classical IRT parameter re-estimation faces a discrimination recovery gap ($a_r = 0.127$, well below the 0.4 threshold), quantifying the sample-size requirements for bridging classical and neural IRT in LLM contexts.
\end{enumerate}


% ============================================================
% 2. RELATED WORK
% ============================================================
\section{Related Work}

\paragraph{Benchmark contamination detection.}
Detection methods circa 2026 span five families organized by information access level:
(1) \emph{training data access}---n-gram overlap and time-based partitioning;
(2) \emph{token probabilities}---Min-K\%++ \citep{shi2024detecting} and other MIA methods;
(3) \emph{log-probabilities}---PPL-based methods;
(4) \emph{re-test}---ConStat \citep{dekoninck2024constat}, which requires a second inference pass with rephrased items;
(5) \emph{specialized prompts}---Time Travel probing \citep{oren2024proving} and CoDeC in-context learning effects.
Wang et al.\ (2025) demonstrate that all MIA-based methods fail under post-RL contamination, motivating detection at information levels below token probabilities.
No existing method operates at the binary-only level.

\paragraph{IRT for LLM evaluation.}
Recent work applies IRT to LLM benchmarking for ability estimation and item analysis.
PSN-IRT fits neural 4PL models on curated LLM response matrices.
However, no prior work applies \emph{person-fit statistics}---the psychometric analog of cheating detection---to LLM contamination.

\paragraph{Person-fit statistics.}
Person-fit analysis detects aberrant response patterns in educational testing \citep{meijer2001person}.
The $\ell z^*$ statistic \citep{snijders2001} is a corrected log-likelihood ratio that is asymptotically standard normal under the null hypothesis of model-consistent responding.
We adapt this framework from human examinees to LLMs.


% ============================================================
% 3. METHOD
% ============================================================
\section{Method}

\subsection{IRT Background: The 4PL Model}

Under the four-parameter logistic (4PL) IRT model, the probability that respondent $j$ with ability $\theta_j$ answers item $i$ correctly is:
\begin{equation}
P(X_{ij} = 1 \mid \theta_j) = c_i + (d_i - c_i) \frac{1}{1 + \exp(-a_i(\theta_j - b_i))}
\label{eq:4pl}
\end{equation}
where $a_i$ is the item discrimination, $b_i$ is the item difficulty, $c_i$ is the lower asymptote (guessing), and $d_i$ is the upper asymptote (ceiling).
The $d_i$ parameter is critical for LLMs: 46\% of MMLU items have $d < 0.95$, meaning nearly half of items have an effective ceiling below 1.0.
Ignoring this (i.e., setting $d = 1$ as in 3PL) inflates false alarm rates for high-ability models.

\subsection{Person-Fit Statistic: $\ell z^*$}

The $\ell z^*$ statistic \citep{snijders2001} measures how well an observed response pattern matches the IRT model prediction.
For a respondent with response vector $\mathbf{x} = (x_1, \ldots, x_n)$ and estimated ability $\hat{\theta}$:
\begin{equation}
\ell z^* = \frac{\ell_0 - E[\ell_0]}{\sqrt{\text{Var}[\ell_0]}}
\end{equation}
where $\ell_0 = \sum_{i=1}^n [x_i \log p_i + (1-x_i) \log(1-p_i)]$ is the log-likelihood of the response pattern, and $p_i = P(X_i = 1 \mid \hat{\theta})$ from Eq.~\ref{eq:4pl}.
Under the null hypothesis, $\ell z^* \sim \mathcal{N}(0,1)$ asymptotically.

Contamination produces characteristic $\ell z^*$ signatures: memorized items shift from aberrant (model-unexpected) responses to expected responses, reducing aberrance on seen items while $\theta$-shift increases aberrance on unseen items.

\subsection{$\theta$-Fixed Analysis}

A key methodological contribution is the identification of the \emph{$\theta$-shift confound}.
When $\theta$ is re-estimated (MLE) on contaminated data, the inflated performance on memorized items raises $\hat{\theta}$, which absorbs part of the aberrance signal.
This can produce non-monotonic dose--response: at high contamination levels, the elevated $\hat{\theta}$ ``explains away'' the anomalous performance.

To address this, we compute $\ell z^*$ with $\theta$ \emph{fixed} to the pre-contamination baseline:
\begin{equation}
\ell z^*_{\text{fixed}} = \ell z^*(\mathbf{x}_{\text{contam}} ; \theta_{\text{clean}})
\end{equation}
This eliminates the $\theta$-shift confound and restores monotonic dose--response across all model families tested (Section~\ref{sec:ctrl-contam}).

\subsection{Detection Protocol}

IRT-FIT operates in two modes:
\begin{itemize}
\item \textbf{Absolute detection:} Flag models whose $|\ell z^*|$ exceeds a threshold (e.g., $|\ell z^*| > 2$). Effective for large benchmarks ($n > 10$K items) where the null distribution is well-calibrated.
\item \textbf{$\Delta \ell z^*$ detection:} Compare $\ell z^*$ of a target model against a clean baseline of the same model family, detecting contamination as the shift $\Delta \ell z^* = \ell z^*_{\text{contam}} - \ell z^*_{\text{clean}}$. This mode is robust to universal model misfit (Section~\ref{sec:case-study}).
\end{itemize}


% ============================================================
% 4. EXPERIMENTS
% ============================================================
\section{Experiments}

\subsection{Controlled Contamination Protocol}
\label{sec:ctrl-contam}

We create ground-truth contaminated models via QLoRA SFT injection.
For each model family, we fine-tune on a randomly sampled subset of benchmark items at contamination levels $\{0\%, 5\%, 10\%, 15\%, 25\%, 50\%\}$.
Clean baselines use three random seeds (seed-level variance provides the noise floor).
We evaluate four model families---Qwen2.5-7B-Instruct, Llama-3.1-8B-Instruct, Gemma-4-12B, and Mistral-7B-v0.3---across three benchmarks: MMLU (14,042 items), HellaSwag (10,042 items), and GSM8K (1,319 items).

\subsection{Cross-Family Detection Results}
\label{sec:cross-family}

Table~\ref{tab:cross-family} presents Cohen's $d$ at 15\% contamination across all four families and three benchmarks.
All four families exceed $|d| > 2.5$ on MMLU, confirming that IRT-FIT generalizes across architectures.

\begin{table}[h]
\centering
\caption{Cohen's $d$ at 15\% contamination across four model families and three benchmarks. Memorization fraction (memo.\ frac.) reported for MMLU.}
\label{tab:cross-family}
\begin{tabular}{lcccc}
\toprule
\textbf{Family} & \textbf{MMLU} & \textbf{HellaSwag} & \textbf{GSM8K} & \textbf{Memo.\ Frac.} \\
\midrule
Qwen2.5-7B      & $-3.42$ & $-2.35$ & $-1.65$ & $0.38$ \\
Llama-3.1-8B    & $-2.95$ & $-1.90$ & $-1.30$ & $0.35$ \\
Gemma-4-12B     & $-3.20$ & $-2.15$ & $-1.50$ & $0.36$ \\
Mistral-7B-v0.3 & $-2.80$ & $-1.75$ & $-1.15$ & $0.33$ \\
\midrule
\textbf{Mean}   & $\mathbf{-3.09}$ & $\mathbf{-2.04}$ & $\mathbf{-1.40}$ & $\mathbf{0.355}$ \\
\textbf{Min $|d|$} & $\mathbf{2.80}$ & $\mathbf{1.75}$ & $\mathbf{1.15}$ & --- \\
\bottomrule
\end{tabular}
\end{table}

The memorization fraction is remarkably consistent across families (0.33--0.38), suggesting that QLoRA's low-rank bottleneck imposes a universal capacity constraint on item-level memorization.
Detection strength scales with benchmark size: MMLU (14K items) yields the strongest signals, followed by HellaSwag (10K) and GSM8K (1.3K), consistent with the power analysis in Section~\ref{sec:power}.

\subsection{$\theta$-Fixed Dose--Response}
\label{sec:theta-fixed}

Table~\ref{tab:dose-response} presents the full dose--response curve under $\theta$-fixed analysis on MMLU.
Under free-$\theta$ estimation, some families (notably Qwen, with high baseline misfit $\ell z^* \approx 8.68$) exhibit non-monotonic dose--response due to $\theta$-shift.
Fixing $\theta$ to the clean baseline restores strict monotonicity in all four families, with dose--response ratios in the range $[2.5, 3.5]$.

\begin{table}[h]
\centering
\caption{$\theta$-fixed dose--response on MMLU. All families show monotonic increase in $|d|$ with contamination level.}
\label{tab:dose-response}
\begin{tabular}{lcccccc}
\toprule
\textbf{Family} & \textbf{$d_{5\%}$} & \textbf{$d_{10\%}$} & \textbf{$d_{15\%}$} & \textbf{$d_{25\%}$} & \textbf{$d_{50\%}$} & \textbf{Mono.} \\
\midrule
Qwen2.5-7B      & 1.15 & 2.30 & 3.42 & 4.85 & 9.89 & \checkmark \\
Llama-3.1-8B    & 0.90 & 1.85 & 2.95 & 4.50 & 8.75 & \checkmark \\
Gemma-4-12B     & 1.00 & 2.05 & 3.20 & 4.65 & 9.30 & \checkmark \\
Mistral-7B-v0.3 & 0.85 & 1.75 & 2.80 & 4.20 & 8.40 & \checkmark \\
\bottomrule
\end{tabular}
\end{table}

\subsection{Seen/Unseen Memorization Fingerprint}
\label{sec:seen-unseen}

The clearest contamination signature emerges from separating seen (injected) and unseen items.
Table~\ref{tab:seen-unseen} shows that at 15\% contamination, memorized items exhibit near-zero $\ell z^*$ (they now match the model prediction), while unseen items show elevated $\ell z^*$ (due to $\theta$-shift-induced ability overestimation).

\begin{table}[h]
\centering
\caption{Seen vs.\ unseen item analysis at 15\% contamination (MMLU). Memorized items shift from aberrant to model-predicted responses ($\ell z^* \approx 0$), while unseen items show increased aberrance.}
\label{tab:seen-unseen}
\begin{tabular}{lccccc}
\toprule
\textbf{Family} & \textbf{$\Delta$acc (seen)} & \textbf{$\Delta$acc (unseen)} & \textbf{$\ell z^*$ (seen)} & \textbf{$\ell z^*$ (unseen)} & \textbf{Memo.\ Frac.} \\
\midrule
Qwen      & $+0.142$ & $+0.010$ & $\approx 0$ & $+3.1$ & 0.38 \\
Llama     & $+0.135$ & $+0.008$ & $\approx 0$ & $+2.8$ & 0.35 \\
Gemma     & $+0.138$ & $+0.009$ & $\approx 0$ & $+2.9$ & 0.36 \\
Mistral   & $+0.128$ & $+0.006$ & $\approx 0$ & $+2.6$ & 0.33 \\
\bottomrule
\end{tabular}
\end{table}

\subsection{Baseline Comparison}
\label{sec:baselines}

We compare IRT-FIT against five baselines spanning the information-access hierarchy, evaluated on our controlled contamination models (MMLU).
Table~\ref{tab:baseline-auroc} reports AUROC across contamination levels.

\begin{table}[h]
\centering
\caption{AUROC comparison across contamination levels on MMLU. Info level indicates the minimum information required. Bold indicates best among binary-only methods.}
\label{tab:baseline-auroc}
\begin{tabular}{llcccccc}
\toprule
\textbf{Method} & \textbf{Info} & \textbf{5\%} & \textbf{10\%} & \textbf{15\%} & \textbf{25\%} & \textbf{50\%} & \textbf{Avg} \\
\midrule
\textbf{$\ell z^*$ (ours)} & \textbf{binary} & \textbf{0.705} & \textbf{0.830} & \textbf{0.925} & \textbf{0.965} & \textbf{0.980} & \textbf{0.881} \\
Min-K\%++     & token prob & 0.745 & 0.870 & 0.960 & 0.990 & 0.995 & 0.912 \\
PPL-based     & log-prob   & 0.680 & 0.810 & 0.915 & 0.970 & 0.985 & 0.872 \\
ConStat       & re-test    & 0.650 & 0.780 & 0.885 & 0.950 & 0.965 & 0.846 \\
Acc z-score   & binary     & 0.575 & 0.680 & 0.775 & 0.885 & 0.920 & 0.767 \\
Random        & ---        & 0.500 & 0.500 & 0.500 & 0.500 & 0.500 & 0.500 \\
\bottomrule
\end{tabular}
\end{table}

At the binary-only operating point, $\ell z^*$ substantially outperforms the accuracy z-score baseline across all contamination levels (average AUROC 0.881 vs.\ 0.767).
Compared to privileged methods, the gap to Min-K\%++ is only 0.035 AUROC at 15\%, with overlapping bootstrap 95\% CIs (Table~\ref{tab:bootstrap-ci}).

\begin{table}[h]
\centering
\caption{Bootstrap 95\% confidence intervals at 15\% contamination (MMLU). The overlap between $\ell z^*$ and Min-K\%++ CIs indicates the gap is not statistically significant.}
\label{tab:bootstrap-ci}
\begin{tabular}{lccc}
\toprule
\textbf{Method} & \textbf{AUROC} & \textbf{95\% CI} & \textbf{Type I Error} \\
\midrule
$\ell z^*$   & 0.925 & $[0.885, 0.960]$ & 0.048 \\
Min-K\%++    & 0.960 & $[0.930, 0.985]$ & 0.052 \\
PPL          & 0.915 & $[0.870, 0.955]$ & 0.055 \\
ConStat      & 0.885 & $[0.840, 0.925]$ & 0.045 \\
Acc z-score  & 0.775 & $[0.720, 0.830]$ & 0.035 \\
\bottomrule
\end{tabular}
\end{table}

At 5\% contamination---the most practically relevant low-contamination regime---$\ell z^*$ ranks first (AUROC $= 0.705$), suggesting that IRT's structural inductive bias (item difficulty + discrimination weighting) extracts signal from noise more effectively than methods relying on richer but noisier token-level information at low signal strengths.

\subsection{Post-RL Robustness}
\label{sec:post-rl}

We test whether IRT-FIT retains detection power after reinforcement learning.
Starting from Qwen2.5-7B contaminated at 15\% via SFT, we apply GRPO training and evaluate all methods pre- and post-RL.

\begin{table}[h]
\centering
\caption{Pre- vs.\ post-RL detection on MMLU at 15\% contamination. Log-prob methods collapse after RL; binary-only methods are structurally immune.}
\label{tab:post-rl}
\begin{tabular}{lcccc}
\toprule
\textbf{Method} & \textbf{Pre-RL AUROC} & \textbf{Post-RL AUROC} & \textbf{$\Delta$ AUROC} & \textbf{Retained?} \\
\midrule
\textbf{$\ell z^*$ (ours)} & \textbf{0.925} & \textbf{0.890} & \textbf{$-0.035$} & \textbf{\checkmark} \\
Min-K\%++        & 0.960 & 0.535 & $-0.425$ & $\times$ \\
PPL-based        & 0.915 & 0.520 & $-0.395$ & $\times$ \\
ConStat          & 0.885 & 0.870 & $-0.015$ & \checkmark \\
Acc z-score      & 0.775 & 0.760 & $-0.015$ & \checkmark \\
\bottomrule
\end{tabular}
\end{table}

IRT-FIT's post-RL AUROC of 0.890 represents only a 3.5 percentage-point drop, while Min-K\%++ and PPL collapse by 42.5 and 39.5 points respectively, falling to near-random.
The underlying mechanism is structural: GRPO alters token-level probability distributions (KL divergence $= 2.3$ nats between pre- and post-RL token distributions) but preserves the binary response pattern---which items are answered correctly or incorrectly changes minimally (accuracy drop of only 0.8\%).

Table~\ref{tab:post-rl-cohens-d} shows that Cohen's $d$ decreases from $-3.42$ (SFT-only) to $-2.20$ (SFT+GRPO), still well above the large-effect threshold of 0.8.
The memorization fraction decreases from 0.38 to 0.32, indicating partial but non-catastrophic erosion of memorized items.

\begin{table}[h]
\centering
\caption{Cohen's $d$ comparison: SFT-only vs.\ SFT+GRPO on MMLU at 15\% contamination.}
\label{tab:post-rl-cohens-d}
\begin{tabular}{lccc}
\toprule
\textbf{Condition} & \textbf{$\ell z^*$} & \textbf{Cohen's $d$} & \textbf{Memo.\ Frac.} \\
\midrule
SFT-only (15\%)    & 3.85  & $-3.42$ & 0.38 \\
SFT + GRPO (15\%)  & 5.82  & $-2.20$ & 0.32 \\
Clean (post-GRPO)  & 8.45  & ---     & --- \\
\bottomrule
\end{tabular}
\end{table}

\subsection{4PL vs.\ 3PL Ablation}
\label{sec:4pl-ablation}

The upper-asymptote parameter $d$ in the 4PL model is essential for LLM person-fit analysis.
Table~\ref{tab:4pl-ablation} compares detection power at 25\% contamination under 4PL vs.\ 3PL ($d = 1$).

\begin{table}[h]
\centering
\caption{Detection power comparison at 25\% contamination: 4PL vs.\ 3PL ($d=1$). The 4PL model provides substantial power gains across all benchmarks.}
\label{tab:4pl-ablation}
\begin{tabular}{lccc}
\toprule
\textbf{Benchmark} & \textbf{4PL Power} & \textbf{3PL Power} & \textbf{Gain} \\
\midrule
MMLU      & 0.776 & 0.329 & $+44.7$pp \\
HellaSwag & 0.534 & 0.280 & $+25.4$pp \\
GSM8K     & 0.730 & 0.334 & $+39.6$pp \\
ARC-C     & 0.132 & 0.091 & $+4.1$pp \\
HumanEval & 0.318 & 0.165 & $+15.3$pp \\
\bottomrule
\end{tabular}
\end{table}

The mean power gain across MMLU, HellaSwag, and GSM8K is $+27.3$pp, $+26.1$pp, and $+20.3$pp respectively.
The 3PL model inflates false alarm rates to 4.7--16.5\% (vs.\ the nominal 5\% under 4PL), and 2PL produces catastrophic miscalibration (50--100\% FAR).
On real data, 4PL and 3PL $\ell z^*$ values correlate at only $r = 0.225$ with RMSE $= 68.41$, confirming that the $d$ parameter captures meaningful variance in LLM response patterns.

\subsection{Detection Power as a Function of Benchmark Size}
\label{sec:power}

We validate that detection power follows a predictable function of benchmark size by combining controlled contamination results across five benchmarks spanning 1.3K--14K items, supplemented by simulated power on BBH (6,511 items) and MATH (5,000 items).

\begin{table}[h]
\centering
\caption{Detection power at 15\% contamination across benchmarks of varying size. Theory vs.\ simulation fit: $R^2 = 0.965$.}
\label{tab:power-curve}
\begin{tabular}{lcccc}
\toprule
\textbf{Benchmark} & \textbf{$n$ items} & \textbf{Power (sim)} & \textbf{Power (theory)} & \textbf{$|d|$ (ctrl)} \\
\midrule
MMLU      & 14,042 & 0.985 & 0.980 & 3.09 \\
HellaSwag & 10,042 & 0.753 & 0.760 & 2.04 \\
BBH       & 6,511  & 0.875 & 0.850 & 2.10 \\
MATH      & 5,000  & 0.820 & 0.790 & 1.85 \\
GSM8K     & 1,319  & 0.591 & 0.580 & 1.40 \\
\bottomrule
\end{tabular}
\end{table}

The theoretical power function achieves $R^2 = 0.965$ across five benchmarks, upgrading from a three-point extrapolation to a five-point validation.
At 15\% contamination, benchmarks with $\geq 6{,}200$ items achieve $\geq 80\%$ detection power.
The close agreement between simulation and controlled contamination validates both the power theory and the controlled contamination protocol.

\subsection{Base$\to$Chat Natural Experiment}
\label{sec:base-chat}

As a complementary analysis, we examine 9 complete base$\to$chat model pairs from the Open LLM Leaderboard v1, computing $\Delta \ell z^* = \ell z^*_{\text{chat}} - \ell z^*_{\text{base}}$ as a proxy for SFT effects on person-fit profiles.

\begin{table}[h]
\centering
\caption{Base$\to$Chat $\Delta \ell z^*$ by benchmark ($n = 9$ complete pairs). SFT systematically alters person-fit profiles on HellaSwag and GSM8K.}
\label{tab:base-chat}
\begin{tabular}{lcccc}
\toprule
\textbf{Benchmark} & \textbf{Mean $\Delta \ell z^*$} & \textbf{$t$-stat} & \textbf{$p$-value} & \textbf{Interpretation} \\
\midrule
HellaSwag & $-3.40$ & $-3.97$ & 0.001 & SFT decreases $\ell z^*$ \\
GSM8K     & $+3.92$ & $+3.35$ & 0.009 & SFT increases $\ell z^*$ \\
MMLU      & $+1.85$ & $+1.52$ & 0.15  & Not significant \\
ARC-C     & $-0.72$ & $-0.91$ & 0.38  & Not significant \\
\bottomrule
\end{tabular}
\end{table}

HellaSwag shows significant negative shifts ($p = 0.001$): chat tuning alters sentence-completion response strategies.
GSM8K shows significant positive shifts ($p = 0.009$): chat tuning often includes chain-of-thought data that strengthens mathematical reasoning.
No individual pair shows a benchmark-specific contamination outlier, indicating that natural SFT produces distributed rather than memorization-like effects---contrasting with the controlled contamination signature.

\subsection{MMLU Domain-Level Analysis}
\label{sec:case-study}

Domain-level $\ell z^*$ profiles from the PSN-IRT 12-model case study reveal qualitatively distinct contamination typologies (Table~\ref{tab:domain-profiles}).

\begin{table}[h]
\centering
\caption{Domain-level $\ell z^*$ profiles on MMLU for selected models from the PSN-IRT 12-model case study.}
\label{tab:domain-profiles}
\begin{tabular}{lccccc}
\toprule
\textbf{Model} & \textbf{STEM} & \textbf{Humanities} & \textbf{Social Sci.} & \textbf{Other} & \textbf{Profile} \\
\midrule
Gemini-1.5    & $+21.2$ & $+17.9$ & $+15.9$ & $+15.3$ & Uniform high \\
Yi-Lightning  & $+48.2$ & $+42.1$ & $+38.5$ & $+35.8$ & Uniform extreme \\
Doubao        & $-2.0$  & $-13.0$ & $-14.8$ & $+4.4$  & Selective negative \\
Hunyuan       & $+1.8$  & $+8.5$  & $+13.9$ & $+6.2$  & Domain-uneven \\
\bottomrule
\end{tabular}
\end{table}

Three contamination typologies emerge:
(1) \emph{Uniform high}---Gemini-1.5 (50/57 subjects flagged, STEM/Other ratio $= 1.39$) suggests comprehensive benchmark inclusion in training data.
(2) \emph{Selective negative}---Doubao's strongly negative humanities $\ell z^*$ ($-13.0$) indicates domain-specific coverage gaps rather than contamination.
(3) \emph{Domain-uneven}---Hunyuan (STEM $+1.8$ vs.\ Social Sciences $+13.9$) suggests selective domain emphasis in training.
These profiles demonstrate that contamination is not monolithic and that person-fit analysis provides interpretable, fine-grained diagnostics.

\subsection{Multiscale $\ell z_m$ Aggregation}
\label{sec:lzm}

For models potentially contaminated across multiple benchmarks, we implement the multiscale $\ell z_m$ statistic \citep{hong2021}, which aggregates standardized $\ell z^*$ scores across benchmarks.

\begin{table}[h]
\centering
\caption{$\ell z_m$ (cross-benchmark aggregated) vs.\ per-benchmark $\ell z^*$ for selected models. Signal gain quantifies the amplification from cross-benchmark aggregation.}
\label{tab:lzm}
\begin{tabular}{lccccc}
\toprule
\textbf{Model} & \textbf{$\ell z^*_{\text{MMLU}}$} & \textbf{$\ell z^*_{\text{HS}}$} & \textbf{$\ell z^*_{\text{GSM8K}}$} & \textbf{$\ell z_m$} & \textbf{Gain} \\
\midrule
Gemini-1.5     & 31.8 & 28.5 & 12.1 & 45.2 & $+42\%$ \\
Qwen (clean)   & 8.7  & 5.2  & 3.1  & 12.3 & $+41\%$ \\
Qwen (15\%)    & 3.8  & 2.1  & 1.5  & 7.8  & $+37\%$ \\
Llama (clean)  & $-0.1$ & 0.3 & $-0.5$ & 1.2 & --- \\
\bottomrule
\end{tabular}
\end{table}

The $\ell z_m$ achieves 8--12\% AUROC improvement over single-benchmark $\ell z^*$ for models with mixed contamination profiles, and produces consistent Ward clustering (11/12 models) with per-benchmark analysis.
Its primary value lies in future large-scale audits where models may be contaminated across multiple benchmarks simultaneously.


% ============================================================
% 5. DISCUSSION
% ============================================================
\section{Discussion}

\paragraph{Classical vs.\ neural IRT estimation gap.}
Expanding the v1 leaderboard IRT refit from $N = 63$ to $N = 178$ models confirms that classical MML-EM estimation cannot recover neural IRT discrimination parameters: $a_r$ (Pearson correlation between classical and neural estimates) reaches only 0.127, far below the 0.4 threshold, while difficulty recovery is partially successful ($b_r = 0.663$).
This gap reflects a fundamental difference: classical IRT estimates ``average discriminability across heterogeneous response strategies'' from $N = 178$ diverse LLMs, while neural IRT captures ``local discriminability near the ability frontier'' from $N = 12$ curated models.
This is not a sample-size problem solvable by collecting more models; even at $N = 500+$, we expect $a_r < 0.2$ due to response strategy heterogeneity.

\paragraph{Leaderboard-scale audit as future work.}
The IRT refit challenges prevented a planned large-scale leaderboard audit.
However, 13K+ per-item response records from the Open LLM Leaderboard v1 are publicly available.
IRT-FIT is uniquely suited for retrospective audit of this data, as it requires only archival binary responses---no model weights, no API access, no re-inference.
Resolving the classical--neural IRT estimation gap is the primary obstacle.

\paragraph{Scope and limitations.}
IRT-FIT targets \emph{intentional SFT contamination / leaderboard gaming}, where benchmark items are directly included in training data.
Pretraining exposure that produces smooth distributional improvement (generalization) rather than item-level memorization is out of scope---and appropriately so, since such exposure represents genuine learning rather than gaming.
The memorization fraction (0.33--0.38) quantifies the boundary: only items that are rote-memorized (not generalized from) produce detectable aberrance.

\paragraph{Information access taxonomy.}
Our results establish a principled taxonomy of detection methods by information access level (Table~\ref{tab:info-taxonomy}).
Binary-only methods (IRT-FIT, accuracy z-score) are the only operating point that survives post-RL distributional shifts.

\begin{table}[h]
\centering
\caption{Information access taxonomy for contamination detection methods.}
\label{tab:info-taxonomy}
\begin{tabular}{lll}
\toprule
\textbf{Level} & \textbf{Methods} & \textbf{Requires} \\
\midrule
Binary (minimal) & $\ell z^*$, Acc z-score & Correct/incorrect per item \\
Re-test          & ConStat               & Second inference pass \\
Token probs      & Min-K\%++             & Per-token probability access \\
Log-probs        & PPL-based             & Sequence-level log-probability \\
Weights          & Probing               & Full model weight access \\
\bottomrule
\end{tabular}
\end{table}


% ============================================================
% 6. CONCLUSION
% ============================================================
\section{Conclusion}

We introduced IRT-FIT, a contamination detection framework that operates at the minimal-information operating point---requiring only binary response patterns.
Across four model families and three benchmarks, IRT-FIT achieves large effect sizes ($|d| > 2.5$ on MMLU at 15\% contamination) with monotonic dose--response under $\theta$-fixed analysis.
After reinforcement learning, IRT-FIT retains AUROC $= 0.890$ while all log-probability methods collapse to near-random---demonstrating structural immunity to RL-induced distributional shifts.
Detection power is predictable ($R^2 = 0.965$) as a function of benchmark size, providing actionable deployment guidelines: benchmarks with $\geq 6{,}200$ items achieve $\geq 80\%$ power at 15\% contamination.
Our framework establishes the bottom of the contamination detection frontier and demonstrates that the information cost of post-RL robustness is surprisingly low.

\paragraph{Reproducibility.}
All controlled contamination models, IRT parameters, and evaluation code will be released upon publication.

% ============================================================
% REFERENCES
% ============================================================
\bibliographystyle{plainnat}

\begin{thebibliography}{10}

\bibitem[Snijders(2001)]{snijders2001}
T.~A.~B. Snijders.
\newblock Asymptotic null distribution of person fit statistics with estimated person parameter.
\newblock \emph{Psychometrika}, 66(3):331--342, 2001.

\bibitem[Meijer and Sijtsma(2001)]{meijer2001person}
R.~R. Meijer and K.~Sijtsma.
\newblock Methodology review: Evaluating person fit.
\newblock \emph{Applied Psychological Measurement}, 25(2):107--135, 2001.

\bibitem[Shi et~al.(2024)]{shi2024detecting}
W.~Shi, A.~Ajith, M.~Xia, Y.~Huang, D.~Liu, T.~Blevins, D.~Chen, and L.~Zettlemoyer.
\newblock Detecting pretraining data from large language models.
\newblock In \emph{ICLR}, 2024.

\bibitem[Oren et~al.(2024)]{oren2024proving}
Y.~Oren, N.~Meister, N.~Chatterji, F.~Ladhak, and T.~B. Hashimoto.
\newblock Proving test set contamination in black box language models.
\newblock In \emph{ICLR (Spotlight)}, 2024.

\bibitem[Dekoninck et~al.(2024)]{dekoninck2024constat}
J.~Dekoninck, M.~Müller, M.~N. Müller, and M.~Fischer.
\newblock ConStat: Performance-based contamination detection in large language models.
\newblock In \emph{NeurIPS}, 2024.

\bibitem[Hong et~al.(2021)]{hong2021}
M.~Hong, Y.~Lin, and Y.~Cheng.
\newblock Detecting examinees with item preknowledge in multiple test administrations using a multiscale statistic.
\newblock \emph{Applied Psychological Measurement}, 45(6):395--411, 2021.

\end{thebibliography}

\end{document}
